\begin{abstract}
\emergencystretch=2em
In multi-agent social settings, model reliability varies across relationships. Beyond inferring what others will do, an agent must calibrate how confidently those inferences should guide policy selection for each relationship. An agent may maintain a well-validated model of one partner, a fragile model of another, and a model under revision for a third; collapsing these into a single confidence estimate loses information relevant to policy selection. We therefore formalize affective precision as a relationship-specific metacognitive estimate of confidence in the current partner model. Each partner's behavioral evidence updates a local confidence estimate that modulates policy precision during selection, regulating how strongly current beliefs are expressed in policy rather than changing the content of those beliefs.

Simulations in a multi-partner graded trust game show that partner-local affective precision influences behavior primarily through policy commitment rather than direct improvement of partner-state inference. Because the mechanism tracks partner-response predictability rather than realized payoff, greater confidence produces sharper policy commitment without necessarily producing higher rewards. Under abrupt shifts in social behavior, confidence accumulated from previously reliable predictions can remain behaviorally active after the relationship changes, revealing a potential cost of stale confidence under social change. Finally, varying precision gain and priors produces distinct trust-calibration dynamics, showing how confidence accumulation and revision depend on model parameters. Together, these results show how relationship-specific affective precision can distinguish social prediction from social policy commitment.

\keywords{Active inference \and Affective precision \and Social metacognition \and Predictive reliability \and Policy precision \and Trust calibration}
\end{abstract}
